% !TeX root = ../main.tex
\section{Simulation}
\label{sec:simulations}

We study two benchmark signals: a varying-smoothness signal (VS), used by \citet{onizuka2024fast}, and a smooth sinusoidal signal (S), used in the first simulation study of \citet{brantley2020baseline}. The former combines gradual variation with a localized feature, whereas the latter represents the smoothly varying baseline that motivates the ADMM approach. Together they assess how Bayesian-PMALA (Section~\ref{subsec:individual}) compares with Bayesian and optimization-based quantile trend filtering under different patterns of smoothness. The interior-quantile experiment uses $n=500$ observations and the five levels $\tau\in\{0.05,0.25,0.50,0.75,0.95\}$. Table~\ref{tab:vs_int_mse} reports VS MSE in the main text; the complete S tables and the MAD and MCIW tables for both signals appear in Appendix~\ref{app:sim:individual}. Section~\ref{sec:sim:extreme} evaluates joint estimation of upper quantiles on both signals, reporting MSE at $\tau\in\{0.96,0.97,0.98,0.99\}$ in Tables~\ref{tab:vs_ext_upper_mse} and \ref{tab:s_ext_upper_mse} and the crossing counts in Appendix~\ref{app:sim:joint_extreme}.

\paragraph{Benchmark signals.} For equally spaced indices $x_i=i$, $i=1,\ldots,n$, let $u_i=i/n$. We generate $y_i=f(u_i)+\varepsilon_i$ independently across indices, with
\begin{align}
 f_{\mathrm{VS}}(u)&=2+\sin(4u-2)+2\exp\{-30(4u-2)^2\},
 \label{eq:sim_vs}\\
 f_{\mathrm{S}}(u)&=\sin(2\pi u).
 \label{eq:sim_s}
\end{align}
We fit both signals with the default order $k=2$ of Section~\ref{subsec:individual}, using the penalty on $D^{(k+1)}\theta$ defined there, so the fitted trend is piecewise quadratic in every configuration. The function $f$ is a baseline location curve; it need not equal the conditional mean or any particular conditional quantile when the errors are asymmetric.

\paragraph{Noise distributions.} Both the ADMM study of \citet{brantley2020baseline} and the Bayesian study of \citet{onizuka2024fast} include Gaussian, Beta, and mixed-normal errors. We use these same three noise families, with the parameterization of \citet{onizuka2024fast}:
\begin{itemize}
 \item Gaussian: $\varepsilon_i\sim\mathcal{N}(0,s(u_i)^2)$, where $s(u)=(1+u^2)/4$.
 \item Beta: $\varepsilon_i\sim\mathrm{Beta}(1,11-10u_i)$.
 \item Mixed normal: $\varepsilon_i\sim u_i\mathcal{N}(-0.2,0.5)+(1-u_i)\mathcal{N}(0.2,0.5)$, where the second argument is the component variance.
\end{itemize}
The Gaussian case is heteroscedastic, the Beta distribution changes shape across the design, and the normal mixture has a design-dependent mixing weight. Brantley et al.\ use the same Gaussian and Beta specifications, but their normal mixture has components $\mathcal{N}(-1,1)$ and $\mathcal{N}(1,1)$; our mixed-normal experiment therefore follows Onizuka et al.\ rather than reproducing that parameterization exactly. All methods in our comparison are evaluated under the same noise specification.

For each level $\tau$, the true conditional quantile is $Q_\tau(u)=f(u)+F_u^{-1}(\tau)$, where $F_u$ is the error distribution at $u$. For Gaussian errors this gives $f(u)+s(u)\Phi^{-1}(\tau)$, and for Beta errors we use the corresponding Beta quantile. For mixed-normal errors we numerically invert
\[
 F_u(t)=u\Phi\!\left(\frac{t+0.2}{\sqrt{0.5}}\right)
       +(1-u)\Phi\!\left(\frac{t-0.2}{\sqrt{0.5}}\right).
\]

\subsection{Interior quantiles}
\label{sec:sim:individual}

\paragraph{Methods and replications.} We compare Bayesian-PMALA with the Laplace and horseshoe MCMC variants of \citet{onizuka2024fast} and the ADMM quantile trend filtering benchmark of \citet{brantley2020baseline}. The point estimates of Bayesian-PMALA and of the two Onizuka samplers are posterior means. Bayesian-PMALA uses the default configuration of Section~\ref{subsec:individual}, with $k=2$ and $\sigma\sim\mathrm{IG}(0.1,0.1)$ in every configuration; the remaining settings are listed in Appendix~\ref{app:pmala}. For the VS signal it uses $\lambda\sim\mathrm{Gamma}(1,10^{-2})$ and $30000$ sweeps of which the first $10000$ are discarded. For the S signal it uses $\lambda\sim\mathrm{Gamma}(1,10^{-3})$, no ASIS step for $\lambda$, and $20000$ sweeps of which the first $5000$ are discarded (Appendix~\ref{app:sim:individual:smooth}). The Onizuka samplers use $8000$ iterations with $2000$ burn-in iterations and their supplied $\mathrm{IG}(1,1)$ prior on $\sigma^2$. ADMM returns an optimization-based point estimate, with $\lambda$ selected by extended BIC over a $25$-point logarithmic grid. All four methods are given the identical simulated response in each replication.

Averages are taken over valid replications. ADMM discards a fit whenever any candidate on its $\lambda$ grid fails to converge, including candidates that the selection criterion does not choose.

\paragraph{Performance measures.} For replication $r$, let $\widehat{\theta}_{\tau,i}^{(r)}$ be the point estimate and $[L_{\tau,i}^{(r)},U_{\tau,i}^{(r)}]$ the $95\%$ pointwise posterior interval. We compute
\begin{align*}
 \mathrm{MSE}^{(r)}_\tau&=\frac{1}{n}\sum_{i=1}^n
  \bigl(\widehat{\theta}_{\tau,i}^{(r)}-Q_\tau(u_i)\bigr)^2,\\
 \mathrm{MAD}^{(r)}_\tau&=\frac{1}{n}\sum_{i=1}^n
  \bigl|\widehat{\theta}_{\tau,i}^{(r)}-Q_\tau(u_i)\bigr|,\\
 \mathrm{MCIW}^{(r)}_\tau&=\frac{1}{n}\sum_{i=1}^n
  \bigl(U_{\tau,i}^{(r)}-L_{\tau,i}^{(r)}\bigr),
\end{align*}
and average each measure over the available replications. MCIW is unavailable for the ADMM point estimator. A smaller interval width alone does not establish better uncertainty quantification; calibration must be assessed separately.

\subsubsection{Varying smoothness}
\label{sec:sim:ind:vs}

Table~\ref{tab:vs_int_mse} reports MSE for the VS signal, and Tables~\ref{tab:vs_int_mad} and \ref{tab:vs_int_mciw} in the appendix report MAD and interval widths. Bayesian-PMALA is comparable in accuracy to the Onizuka samplers across noise families and quantile levels, although it uses a single global penalty where the horseshoe assigns a local shrinkage scale to every difference, and it is more accurate than ADMM. The localized feature of the VS signal calls for a penalty light enough to resolve it; the posterior learns that penalty from the data jointly with the likelihood scale, whereas extended BIC selects a heavier one.

\begin{table}[!htbp]
\centering
\small
\caption{Mean squared error (MSE) for the varying-smoothness (VS) signal ($k=2$), $n=500$, at interior quantile levels.}
\label{tab:vs_int_mse}
\begin{tabular}{lrrrrr}
\toprule
Beta & $0.05$ & $0.25$ & $0.50$ & $0.75$ & $0.95$ \\
\midrule
Onizuka (Laplace) & 0.000513 & 0.001024 & 0.001942 & 0.002994 & 0.007269 \\
Onizuka (Horseshoe) & 0.000512 & 0.000943 & 0.001802 & 0.002778 & 0.006493 \\
Bayesian-PMALA & 0.000546 & 0.001035 & 0.001930 & 0.003029 & 0.007953 \\
ADMM & 0.008055 & 0.004120 & 0.006632 & 0.007613 & 0.014649 \\
\midrule
Mixed normal & $0.05$ & $0.25$ & $0.50$ & $0.75$ & $0.95$ \\
\midrule
Onizuka (Laplace) & 0.075391 & 0.031684 & 0.027303 & 0.031800 & 0.077075 \\
Onizuka (Horseshoe) & 0.073529 & 0.031743 & 0.028090 & 0.033608 & 0.079668 \\
Bayesian-PMALA & 0.079580 & 0.033127 & 0.026996 & 0.031119 & 0.074788 \\
ADMM & 0.096462 & 0.092891 & 0.071678 & 0.049944 & 0.068415 \\
\midrule
Gaussian & $0.05$ & $0.25$ & $0.50$ & $0.75$ & $0.95$ \\
\midrule
Onizuka (Laplace) & 0.020916 & 0.008284 & 0.006932 & 0.007558 & 0.018128 \\
Onizuka (Horseshoe) & 0.019990 & 0.008067 & 0.006651 & 0.007358 & 0.017762 \\
Bayesian-PMALA & 0.020012 & 0.007919 & 0.006843 & 0.007612 & 0.018994 \\
ADMM & 0.088108 & 0.025091 & 0.020181 & 0.014800 & 0.022334 \\
\bottomrule
\end{tabular}
\end{table}

\clearpage
\subsubsection{Smooth signal}
\label{sec:sim:ind:smooth}

Results for the S signal are in Appendix Tables~\ref{tab:s_int_mse}, \ref{tab:s_int_mad} and \ref{tab:s_int_mciw}. ADMM is the most accurate method on this signal, as expected for the smooth-baseline setting for which it was designed \citep{brantley2020baseline}, and Bayesian-PMALA is more accurate than both Onizuka samplers in most configurations: without a localized feature the local shrinkage scales have little to exploit, and one global penalty learned from the data suffices.

Across the two signals each competitor is most accurate on the benchmark that suits it, the horseshoe on VS and ADMM on S, while Bayesian-PMALA is competitive on both with one default configuration. It matches the locally adaptive sampler with a simpler hierarchy and banded computations (Section~\ref{sec:scalability}), and it matches the optimizer while returning posterior draws of the trend, the penalty and the likelihood scale instead of a point estimate at a selected penalty. The empirical coverage of its pointwise $95\%$ intervals is $0.86$--$0.94$ at $\tau=0.25$, $0.50$ and $0.75$, and $0.65$--$0.72$ at $\tau=0.05$ and $0.95$ except under Beta noise at $\tau=0.05$, where it is $0.94$.

\subsection{Joint estimation of extreme upper quantiles}
\label{sec:sim:extreme}

We use the joint sampler of Section~\ref{subsec:joint} to assess accuracy at extreme upper quantiles and to exploit its non-crossing constraint when quantile levels are close together. For both signals and the three noise distributions above, we jointly fit $\tau\in\{0.90,0.93,0.96,0.97,0.98,0.99\}$ and evaluate the four levels $0.96$--$0.99$; the two lower levels are fitted but not evaluated. The comparison includes Bayesian-PMALA, the Laplace and horseshoe samplers of \citet{onizuka2024fast} and the ADMM estimator of \citet{brantley2020baseline}, each fitted one level at a time, and all five methods are given the identical simulated response in each replication. Tables~\ref{tab:vs_ext_upper_mse} and \ref{tab:s_ext_upper_mse} report MSE for the VS and S signals, and Appendix~\ref{app:sim:joint_extreme} reports how often the fitted curves cross.

On the VS signal the joint sampler has the lowest MSE in $9$ of the $12$ configurations: all Gaussian and mixed-normal cases and the Beta case at $\tau=0.99$, where it and individual Bayesian-PMALA are nearly equal. The horseshoe sampler leads under Beta noise at $\tau=0.96$--$0.98$, as it does under Beta noise at interior levels (Section~\ref{sec:sim:ind:vs}). The margin over the best competitor grows with $\tau$, from $11\%$ at $\tau=0.96$ to $26\%$ at $\tau=0.99$ under Gaussian noise and from $18\%$ to $42\%$ under mixed-normal noise.

On the S signal ADMM is the most accurate method under Beta and Gaussian noise at every level and under mixed-normal noise at $\tau=0.96$, as at interior levels (Section~\ref{sec:sim:ind:smooth}); the joint sampler is second in each of these cases and nearly equal to ADMM under Gaussian noise at $\tau=0.99$. Under mixed-normal noise the ADMM error rises from $0.035$ at $\tau=0.96$ to $0.291$ at $\tau=0.99$, against $0.043$ to $0.094$ for the joint sampler, which has the lowest MSE at $\tau=0.97$--$0.99$. On this signal the joint sampler is more accurate than the three Bayesian competitors in every configuration.

The joint sampler improves on individual Bayesian-PMALA in all $24$ configurations. The two runs differ in length and tuning as well as in the constraint, so the comparison does not isolate the effect of the ordering. The gains come with ordered fits: the joint posterior means do not cross at any index, whereas every individually fitted method produces crossing curves at these closely spaced levels in every configuration, the horseshoe sampler least often and Bayesian-PMALA or ADMM most often (Appendix~\ref{app:sim:joint_extreme}).

\begin{table}[!htbp]
\centering
\small
\caption{Mean squared error at extreme upper quantiles for the VS signal, $n=500$ and $k=2$. The joint sampler fits six levels together.}
\label{tab:vs_ext_upper_mse}
\begin{tabular}{lrrrr}
\toprule
Beta & $0.96$ & $0.97$ & $0.98$ & $0.99$ \\
\midrule
Onizuka (Laplace) & 0.008704 & 0.010190 & 0.011954 & 0.015642 \\
Onizuka (Horseshoe) & 0.007479 & 0.008790 & 0.010444 & 0.013843 \\
Bayesian-PMALA & 0.010588 & 0.012392 & 0.015399 & 0.013241 \\
ADMM & 0.016921 & 0.017415 & 0.019852 & 0.018345 \\
Joint sampler (non-crossing) & 0.008336 & 0.011275 & 0.012640 & 0.013170 \\
\midrule
Mixed normal & $0.96$ & $0.97$ & $0.98$ & $0.99$ \\
\midrule
Onizuka (Laplace) & 0.099213 & 0.131297 & 0.169168 & 0.243745 \\
Onizuka (Horseshoe) & 0.093988 & 0.118348 & 0.144101 & 0.212250 \\
Bayesian-PMALA & 0.120335 & 0.146784 & 0.184038 & 0.245757 \\
ADMM & 0.092590 & 0.124361 & 0.184480 & 0.298003 \\
Joint sampler (non-crossing) & 0.076015 & 0.087959 & 0.103372 & 0.124001 \\
\midrule
Gaussian & $0.96$ & $0.97$ & $0.98$ & $0.99$ \\
\midrule
Onizuka (Laplace) & 0.022541 & 0.028454 & 0.034980 & 0.049972 \\
Onizuka (Horseshoe) & 0.021205 & 0.025912 & 0.031785 & 0.045252 \\
Bayesian-PMALA & 0.027103 & 0.032784 & 0.040975 & 0.050202 \\
ADMM & 0.023359 & 0.026523 & 0.032114 & 0.040680 \\
Joint sampler (non-crossing) & 0.018820 & 0.022292 & 0.026460 & 0.029993 \\
\bottomrule
\end{tabular}
\end{table}

\begin{table}[!htbp]
\centering
\small
\caption{Mean squared error at extreme upper quantiles for the S signal, $n=500$ and $k=2$. The joint sampler fits six levels together.}
\label{tab:s_ext_upper_mse}
\begin{tabular}{lrrrr}
\toprule
Beta & $0.96$ & $0.97$ & $0.98$ & $0.99$ \\
\midrule
Onizuka (Laplace) & 0.006104 & 0.007880 & 0.009660 & 0.012155 \\
Onizuka (Horseshoe) & 0.006160 & 0.007535 & 0.009146 & 0.011910 \\
Bayesian-PMALA & 0.008803 & 0.010699 & 0.013438 & 0.012056 \\
ADMM & 0.003221 & 0.003291 & 0.003851 & 0.004996 \\
Joint sampler (non-crossing) & 0.003794 & 0.004098 & 0.004853 & 0.006763 \\
\midrule
Mixed normal & $0.96$ & $0.97$ & $0.98$ & $0.99$ \\
\midrule
Onizuka (Laplace) & 0.090608 & 0.125382 & 0.165706 & 0.240934 \\
Onizuka (Horseshoe) & 0.085961 & 0.110166 & 0.138938 & 0.200308 \\
Bayesian-PMALA & 0.111903 & 0.141999 & 0.176497 & 0.232418 \\
ADMM & 0.034574 & 0.071379 & 0.173204 & 0.291077 \\
Joint sampler (non-crossing) & 0.043013 & 0.051190 & 0.064297 & 0.094006 \\
\midrule
Gaussian & $0.96$ & $0.97$ & $0.98$ & $0.99$ \\
\midrule
Onizuka (Laplace) & 0.017660 & 0.023910 & 0.031704 & 0.043506 \\
Onizuka (Horseshoe) & 0.017837 & 0.022887 & 0.028483 & 0.040580 \\
Bayesian-PMALA & 0.023762 & 0.029324 & 0.037842 & 0.045966 \\
ADMM & 0.007117 & 0.008640 & 0.011142 & 0.021601 \\
Joint sampler (non-crossing) & 0.010749 & 0.012057 & 0.015037 & 0.021862 \\
\bottomrule
\end{tabular}
\end{table}
